% !TeX root = ../main.tex
\section{Conclusions and outlook}
\label{sec:conclusions}

We developed a nonequilibrium, finite-deformation theory for
multicomponent, multiphase porous media with reactions, phase
transformations, two thermal subsystems, and electrostatics.  The central
contribution is a traceable division of the model into conservation laws,
reversible constraints, thermodynamic restrictions, and dissipative closures.
The constrained Hamilton--d'Alembert statement produces the phase-coordinate,
momentum, storage, composition, transfer-potential, and electrostatic
relations, including the d'Alembert work of the saturation resistance when
that branch is active.  The Coleman--Noll procedure identifies the material
stresses, phase pressures, entropies, generalized component-transfer forces,
electrochemical affinities, gauge-invariant nonisothermal transport forces,
capillary-history dissipation, and residual dissipation.  The entropy
inequality supplies the admissibility and nonnegative-dissipation restrictions.
Onsager stationarity selects near-equilibrium rates for dynamic saturation
redistribution, reaction, diffusion, and dispersion.  The independently
derived reversible relations agree, which supplies an internal check on the
construction.

Converted mass enters the phase momentum production and the generalized Darcy
law.  The resulting source-shifted resistance distinguishes dissipation of the
underlying drag from algebraic nonsingularity and the stronger
positive-resistance regime.  For the solids, the fixed-pressure Legendre
construction gives nonlinear Biot coefficients at the current reacted state,
while the phase/component kinematics connect pore-filling conversion to
distension, stress-free deformation, elasticity, and plastic flow.  The
true-plastic and distension powers identify two conjugate stress--rate pairs and
provide the basis for a coupled generalized-plasticity specialization.

The thermal, chemical, and electrical results use the same component-transfer
and electric-enthalpy structure.  Composition stationarity recovers
electrochemical equilibrium and supplies the transfer-work correction away
from equilibrium.  Separating the nonisothermal component-potential force from
the electric field preserves gauge invariance.  The constrained interfacial
energy gives the stored capillary pressure and capillary-history dissipation;
the optional saturation resistance adds a dynamic pressure lag without changing
the stored energy.

The construction is internally consistent.  The common charge-rate and
electrostatic-power identities couple electrical body force and source power
without double counting.  The same phase electric-enthalpy family supplies
the electrocapillary and momentum relations, and the component material
measures control both conversion momentum and solid compositional evolution.
Writing the energy and entropy balances in a skeleton observer retains bulk
phase transport relative to the solid and exposes its joint admissibility
condition with mechanical interaction-energy exchange.

The specialized reductions check the general equations.  Under their stated
assumptions, the theory recovers the standard mass-based compositional
balance, Fick--Onsager transport in \(N-1\) independent composition
directions, black-oil conservation equations and Darcy fluxes, a
Nernst--Planck--Darcy system, and the finite-deformation and linearized
single-solid/single-fluid Biot equations.  The neutral one-fluid limit removes
the electrocapillary and interfacial terms, and the local-thermal-equilibrium
limit recovers the usual affinity power for phasewise mass-conserving
chemistry or vanishing transfer-work offsets.  These recoveries identify the
thermodynamic structure retained when the full model is specialized.

Several limitations delimit the present result.  The theory applies while the
participating phases and components remain present; phase appearance and
disappearance require separate switching conditions.  Capillary hysteresis is
represented by continuous internal variables and material-specific
loading--unloading rules.  The dynamic pressure lag is an optional
near-equilibrium saturation-rate resistance whose calibration determines
whether the branch is active for a given reservoir transient.  Independent
interfacial-area evolution requires a further state variable and kinetic
equation.  The electrical state is electroquasistatic and isotropic and
excludes permanent polarization, directional dielectric structure,
relaxation, and wave dynamics.  Inserted mass carries recipient-phase energy
and entropy under a local-equilibrium source convention.  Reactions spanning
the two thermal subsystems use the gauge-invariant neutral
temperature-weighted reaction force; their physical reaction-energy transfer
may be Onsager-coupled under the stated mechanismwise dissipation condition.
The quadratic closures shown here are limited to near-equilibrium response.
The finite-deformation Biot storage tangent used in the classical comparison
is an additional integrable storage assumption introduced independently of
the general phase energy.  Constitutive calibration, well-posedness,
discretization, and benchmark validation remain outside the present
derivation.

Implementation should proceed with equation-level traceability.  A
solid-reference discretization can first verify neutral-isothermal component
balances, relative-flux reconstruction, and coupled momentum before adding
reaction-induced solid-volume change, two-temperature exchange, and
electrochemical transport.  Manufactured solutions and the derived limits then
provide regression targets for pore-filling, crystallization, and
electrocapillary benchmarks.  A cyclic two-fluid drainage--imbibition test can
verify reversal continuity, scanning paths, trapping, and nonnegative
capillary-history dissipation before the closure is extended to three phases.
A rate sweep for fracture--matrix imbibition and flowback in a tight or shale
sample can then compare \(\mc T=\mbf0\) with the dynamic closure using uptake,
pressure lag, and saturation profiles as observables.
